\chapter{Equations of Mean Curvature Type}

The prescribed mean-curvature equation is
\begin{equation}
\mathfrak Mu=(1+|Du|^2)\Delta u-D_iuD_juD_{ij}u
=nH(1+|Du|^2)^{3/2}.
\tag{16.1}
\end{equation}


\section*{16.1. Hypersurfaces in $\mathbb R^{n+1}$}

\begin{definition}[Tangential differential operators]\label{gilbarg-trudinger.def.tangential-gradient}\dcref{ch16:16.2--16.4}\group{chapter-16}
Let $\mathcal U\subset\mathbb R^{n+1}$ be open and let
$\Phi\in C^2(\mathcal U)$ have nonvanishing gradient.  On
\begin{equation}
\mathcal S=\{x\in\mathcal U:\Phi(x)=0\},
\tag{16.2}
\end{equation}
take the unit normal $\nu=D\Phi/|D\Phi|$.  A graph is represented by
\begin{equation}
\Phi(x',x_{n+1})=x_{n+1}-u(x'),
\qquad x'\in\Omega.
\tag{16.3}
\end{equation}
For $g\in C^1(\mathcal U)$, define
\begin{equation}
\delta g=Dg-(\nu\cdot Dg)\nu.
\tag{16.4}
\end{equation}
\end{definition}

This is the tangential projection of $Dg$, depends only on
$g|_{\mathcal S}$, and satisfies
\begin{equation}
\nu\cdot\delta g=0,
\tag{16.5}
\end{equation}
\begin{equation}
|\delta g|^2=|Dg|^2-|\nu\cdot Dg|^2,
\tag{16.6}
\end{equation}
\begin{equation}
|\delta g|\le|Dg|,
\tag{16.7}
\end{equation}
\begin{equation}
|\nu_{n+1}|\,|Dg|\le|\delta g|
\quad\text{if }D_{n+1}g=0.
\tag{16.8}
\end{equation}
The mean curvature with respect to $\nu$ is
\begin{equation}
H=-\frac1n\delta_i\nu_i.
\tag{16.9}
\end{equation}

\begin{lemma}[Tangential integration by parts]\label{gilbarg-trudinger.lem.16.1.tangential-integration-by-parts}\dcref{ch16:16.1}\group{chapter-16}
For $g\in C^1_0(\mathcal U)$ and area element $dA$ on $\mathcal G$,
\begin{equation}\tag{16.10}
\int_{\mathcal G}\delta g\,dA=-n\int_{\mathcal G}gH\nu\,dA.
\end{equation}
\end{lemma}
\begin{proof}
\sketch On a graph, set $v=\sqrt{1+|Du|^2}$ and use $\nu_i=-D_i u/v$, $dA=v\,dx$, and
\begin{equation}\tag{16.11}
\nu_{n+1}=\frac1v.
\end{equation}
Integration by parts gives $\int_{\mathcal G}\delta_i g\,dA=-n\int_{\mathcal G}gH\nu_i\,dA$ componentwise; a partition of unity gives the general case.
\end{proof}

Note that the case $i=n+1$ in Lemma 16.1 is equivalent to the integral form of
the prescribed mean curvature equation.
\begin{definition}[Laplace--Beltrami operator on a hypersurface]\label{gilbarg-trudinger.def.laplace-beltrami}\dcref{ch16:16.12}\group{chapter-16}
For $g\in C^2(\mathcal U)$ define
\begin{equation}\tag{16.12}
\Delta g=\delta_i\delta_i g.
\end{equation}
\end{definition}

For $\varphi\in C_0^2(\mathcal U)$, (16.10) gives
\begin{equation}
\int_{\mathcal G}\varphi\Delta g\,dA
=\int_{\mathcal G}g\Delta\varphi\,dA
=-\int_{\mathcal G}\delta g\cdot\delta\varphi\,dA.
\tag{16.13}
\end{equation}
For $y\in\mathcal G$, $r=|x-y|$, and a radial function $\psi(r)$,
\begin{equation}
\Delta\psi(r)=\frac{n\psi'}r
+\left(\psi''-\frac{\psi'}r\right)|\delta r|^2
+\frac{n\psi'}rH\nu\cdot(x-y),
\tag{16.14}
\end{equation}
where $|\delta r|^2=1-|\nu\cdot(x-y)|^2/r^2$.  For nonnegative,
nonincreasing $\chi\in C^1(\mathbb R)$ supported in $(-\infty,1)$, set
\[
\psi(r)=\int_r^\infty\tau\chi(\tau/\rho)\,d\tau,
\]
where $0<\rho<R$ and $B_R(y)\subset\mathcal U$.  Then
\begin{equation}
\begin{aligned}
\Delta\psi(r)
&=\rho^{n+1}D_\rho[\rho^{-n}\chi(r/\rho)]
+\frac r\rho\chi'(r/\rho)(1-|\delta r|^2)\\
&\qquad-n\chi(r/\rho)H\nu\cdot(x-y)\\
&\le\rho^{n+1}D_\rho[\rho^{-n}\chi(r/\rho)]
-n\chi(r/\rho)H\nu\cdot(x-y).
\end{aligned}
\tag{16.15}
\end{equation}
When $H=0$ this reduces to
\begin{equation}
\Delta\psi(r)\le\rho^{n+1}D_\rho[\rho^{-n}\chi(r/\rho)].
\tag{16.16}
\end{equation}
Suppose $g\ge0$ is weakly subharmonic on a minimal hypersurface, so
\begin{equation}
\int_\mathcal Gg\Delta\varphi\,dA\ge0
\tag{16.17}
\end{equation}
for nonnegative $\varphi\in C_0^2(\mathcal U)$.  Equations (16.16)--(16.17)
make
\begin{equation}
I_\chi(\rho)=\frac1{\omega_n\rho^n}
\int_\mathcal Sg\chi(r/\rho)\,dA
\tag{16.18}
\end{equation}
nondecreasing.  Approximation by characteristic cutoffs gives the same property
for
\begin{equation}
I(\rho)=\frac1{\omega_n\rho^n}
\int_{\mathcal S\cap B_\rho(y)}g\,dA,
\tag{16.19}
\end{equation}
and differentiation of surface measure gives
\begin{equation}
\lim_{\rho\to0}I(\rho)=g(y).
\tag{16.20}
\end{equation}
Consequently
\begin{equation}
g(y)\le\frac1{\omega_nR^n}
\int_{\mathcal S\cap B_R(y)}g\,dA.
\tag{16.21}
\end{equation}

\begin{lemma}[Mean-value inequality for minimal hypersurfaces]\label{gilbarg-trudinger.lem.16.2.mean-value}\dcref{ch16:16.2}\group{chapter-16}
Let $g\geq0$ be subharmonic ($\Delta g\geq0$) on a $C^2$ minimal hypersurface $\mathcal S$. If $y\in\mathcal S$ and $B_R(y)\subset\mathcal U$, inequality (16.21) holds for almost every $y$ (and every $y$ when $g$ is continuous).
\end{lemma}
\begin{proof}
\sketch The cutoff computation (16.15)--(16.16) makes $I_\chi$ in (16.18) monotone. Approximation by characteristic cutoffs and the limit (16.20) yield (16.21).
\end{proof}

Taking $g=1$ in (16.21) gives
\begin{equation}
A(\mathcal S\cap B_R(y))\ge\omega_nR^n.
\tag{16.22}
\end{equation}
Write $\mathcal S_R(y)=\mathcal S\cap B_R(y)$.  For a general
hypersurface, let $\chi\in C^1[0,\infty)$ be nonnegative, nonincreasing,
and supported in $[0,R]$, and set
\[
\psi(r)=\int_r^\infty \tau\chi(\tau)\,d\tau,
\]
Then (16.15) becomes
\begin{equation}\tag{16.23}
\Delta\psi(r)= -\{n\chi(r)+r\chi'(r)|\delta r|^2+n\chi(r)Hv\cdot (x-y)\}
=-(n\chi(r)+r\chi'(r))+r\chi'(r)(1-|\delta r|^2)-n\chi(r)Hv\cdot(x-y).
\end{equation}

Therefore, substituting into (16.13), we obtain

\begin{equation}\tag{16.24}
\int_{\Scal}\bigl(n\chi(r)+r\chi'(r)\bigr)g\,dA-\int_{\Scal}r\chi'(r)(1-|\delta r|^2)g\,dA
=-\,n\int_{\Scal}\chi(r)gHv\cdot(x-y)\,dA+\int_{\Scal}\delta\psi\cdot\delta g\,dA.
\end{equation}

If $0<\varepsilon<R$ and
$\chi=\varepsilon^{-n}-R^{-n}$ on $[0,\varepsilon)$, then
\begin{equation}\tag{16.25}
n(\varepsilon^{-n}-R^{-n})\int_{\Scal_\varepsilon}g\,dA+\int_{\Scal_R-\Scal_\varepsilon}\bigl(n\chi(r)+r\chi'(r)\bigr)g\,dA
-\int_{\Scal_R-\Scal_\varepsilon}r\chi'(r)(1-|\delta r|^2)g\,dA
=-\,n\int_{\Scal_R}\chi(r)gHv\cdot(x-y)\,dA+\int_{\Scal_R}\delta\psi\cdot\delta g\,dA.
\end{equation}


Choose the limiting cutoff
\[
\chi_\varepsilon(r)=
\begin{cases}
\varepsilon^{-n}-R^{-n},&0\le r<\varepsilon,\\
r^{-n}-R^{-n},&\varepsilon\le r<R,\\
0,&r\ge R.
\end{cases}
\]
Smooth monotone approximation in (16.25) yields

\[
(\varepsilon^{-n}-R^{-n})\int_{\Ss_\varepsilon} g\,dA+\int_{\Ss_R-\Ss_\varepsilon} gr^{-n}(1-|\delta r|^2)\,dA
\]
\begin{equation}
=R^{-n}\int_{\Ss_R-\Ss_\varepsilon} g\,dA-(\varepsilon^{-n}-R^{-n})\int_{\Ss_\varepsilon} g\Hv\cdot(x-y)\,dA
\notag\\
-\int_{\Ss_R-\Ss_\varepsilon} g(r^{-n}-R^{-n})\Hv\cdot(x-y)\,dA+\frac{1}{n}\int_{\Ss_R}\delta\psi_\varepsilon(r)\cdot\delta g\,dA,
\tag{16.26}
\end{equation}

where

\[
\psi_\varepsilon(r)=\int_r^R t\chi_\varepsilon(t)\,dt.
\]

Letting $\varepsilon\downarrow0$ gives

\[
g(y)+\frac{1}{\omega_n}\int_{\Ss_R}\frac{(1-|\delta r|^2)}{r^n}\,g\,dA
\]
\begin{equation}
=\frac{1}{\omega_n R^n}\int_{\Ss_R} g\,dA-\frac{1}{\omega_n}\int_{\Ss_R} g(r^{-n}-R^{-n})\Hv\cdot(x-y)\,dA
\notag\\
+\frac{1}{n\omega_n}\int_{\Ss_R}\delta\psi(r)\cdot\delta g\,dA,
\tag{16.27}
\end{equation}

where

\[
\psi(r)=\int_r^R \tau\bigl(\tau^{-n}-R^{-n}\bigr)\,d\tau.
\]

Using

\begin{equation}
|Hv\cdot (x-y)|\le r^{-2}|v\cdot (x-y)|^2+\frac14 H^2r^2
=1-|\delta r|^2+\frac14 H^2r^2
\tag{16.28}
\end{equation}

in (16.27) gives

\begin{equation}
g(y)\le \frac{1}{\omega_n R^n}\int_{\mathcal{G}_R} g|\delta r|^2\,dA
+\frac{1}{4\omega_n}\int_{\mathcal{G}_R} gH^2r^2(r^{-n}-R^{-n})\,dA
-\frac{1}{n\omega_n}\int_{\mathcal{G}_R}(r^{-n}-R^{-n})(x-y)\cdot \delta g\,dA.
\tag{16.29}
\end{equation}

\begin{lemma}[Potential estimate on a hypersurface]\label{gilbarg-trudinger.lem.16.3.potential-estimate}\dcref{ch16:16.3}\group{chapter-16}
Let $g\in C^1(\mathcal U)$ be nonnegative. For every $y\in\mathcal G$ and $B_R(y)\subset\mathcal U$, inequality (16.29) holds.
\end{lemma}
\begin{proof}
\sketch Insert the radial cutoff $\chi_\varepsilon(r)=\varepsilon^{-n}-R^{-n}$ for $r<\varepsilon$ and $r^{-n}-R^{-n}$ for $\varepsilon\leq r<R$ into (16.24). Passing to the limit gives (16.27); estimate the curvature term by (16.28) to obtain (16.29).
\end{proof}

For $g\in C^2$, the last term in (16.29) equals
$-(n\omega_n)^{-1}\int_{\mathcal G_R}\psi\Delta g\,dA$.
In dimension two, (16.29) gives
\begin{equation}
g(y)\le \frac{1}{\pi R^2}\int_{\mathcal{G}_R} g\left(1+\frac14 H^2R^2\right)\,dA
+\frac{1}{2\pi}\int_{\mathcal{G}_R} r(r^{-2}-R^{-2})|\delta g|\,dA.
\tag{16.30}
\end{equation}

For $g=1$,
\begin{equation}
1\le \frac{A(\mathcal{G}_R)}{\pi R^2}+\frac{1}{4\pi}\int_{\mathcal{G}_R} H^2\,dA.
\tag{16.31}
\end{equation}

Thus a compact hypersurface in $\mathbb R^3$, or one with
$A(\mathcal S_R)=o(R^2)$, satisfies
\begin{equation}
\int_{\mathcal{S}} H^2\,dA \geq 4\pi.
\tag{16.32}
\end{equation}

with equality precisely for a sphere.  From (16.26), (16.28), and Young's
inequality,

\begin{equation}
\tag{16.33}
\begin{aligned}
\sup_{(0,R)} I(\rho)&\leq \frac{1}{\omega_n R^n}\int_{\mathcal{S}_R}\big[g|n\delta r|^2+|HR|^n+nH^2r^2R^n(r^{-n}-R^{-n})\big]\,dA\\
&\quad+\frac{1}{\omega_n}\int_{\mathcal{S}_R} r(r^{-n}-R^{-n})|\delta g|\,dA.
\end{aligned}
\end{equation}

For $g\in C^2$, the last term may be replaced by
$\omega_n^{-1}\int_{\mathcal S_R}\psi(-\Delta g)^+\,dA$.
For $n=2$ and $g=1$ this yields

\begin{equation}
\sup_{(0,R)} \frac{A(\mathcal{S}_\rho)}{\rho^2}\leq \frac{2A(\mathcal{S}_R)}{R^2}+3\int_{\mathcal{S}_R} H^2\,dA.
\tag{16.34}
\end{equation}

Set
\[
J(\rho)=D_\rho\left[\int_{\mathcal{S}_\rho} g|\delta r|^2\,dA\right].
\]
The same cutoff argument gives

\begin{equation}
J(\rho)\le \frac{1}{\rho}\int_{\mathcal S_{\rho}}\bigl[n|g|(1+|H|r)+r|\delta g|\bigr]\,dA
\tag{16.35}
\end{equation}

\[
\le \frac{n}{\rho}\int_{\mathcal S_{\rho}}|g|\,dA+\int_{\mathcal S_{\rho}}(n|Hg|+|\delta g|)\,dA.
\]

For $g=1$ and $n=2$, (16.34)--(16.35) imply
\begin{equation}
D_\rho\int_{\mathcal S_\rho}|\delta r|^2\,dA
\le10\rho\left\{\frac{A(\mathcal G_\rho)}{R^2}
+\int_{\mathcal S_R}H^2\,dA\right\}.
\tag{16.36}
\end{equation}

The following surface Morrey estimate is used for quasiconformal maps.

\begin{lemma}[Morrey estimate on a surface]\label{gilbarg-trudinger.lem.16.4.surface-morrey-estimate}\dcref{ch16:16.4}\group{chapter-16}
Let $g\in C^{1}(\mathcal U)$, $n=2$, and suppose there are $K>0$ and $\beta\in(0,1)$ such that

\begin{equation}
\int_{\mathcal S_{\rho}(y)}|\delta g|\,dA\le K\rho(\rho/R)^{\beta}
\tag{16.37}
\end{equation}

for every $\bar y\in \mathcal G_{R/4}(y)$ and $\rho\le R/4$. Then

\begin{equation}
\sup_{x\in \mathcal S^*_{\rho}(y)}|g(x)-g(y)|\le CK(\rho/R)^{\beta}\left\{\frac{A(\mathcal G_{R})}{R^{2}}+\int_{\mathcal S_{R}}H^{2}\,dA\right\}
\tag{16.38}
\end{equation}


where $C=C(\beta)$ and $\mathcal{S}^*_\rho(y)$ is the component of $\mathcal{S}_\rho(y)$ containing $y$.
\end{lemma}

\begin{proof}
\sketch Rewrite (16.30) as

\begin{equation}
g(y)\le \frac{1}{\pi}\left\{\int_{\mathcal{S}_R} g\left(R^{-2}+H^2/4\right)\,dA
+ \int_0^R \rho^{-2}\int_{\mathcal{S}_\rho} |\delta g|\,dA\,d\rho\right\},
\tag{16.39}
\end{equation}

Partition the oscillation of $g$ on $\mathcal S^*_\rho(y)$ into intervals of length
$6\beta^{-1}K(\rho/R)^\beta$. Applying (16.39) to smooth cutoffs of these intervals and
summing, with (16.34), bounds their number by the right-hand side of (16.38), proving the estimate.

\end{proof}

\begin{remark}[Compact-support potential identities]\label{gilbarg-trudinger.remark.compact-support-potential}\dcref{ch16:16.40}\group{chapter-16}
If $g$ has compact support, the limit $R\to\infty$ in (16.27)--(16.29)
gives
\begin{equation}
\begin{aligned}
g(y)+\frac1{\wn}\int_\mathcal S
\frac{1-|\delta r|^2}{r^n}g\,dA
&=\frac1{\wn}\int_\mathcal S
\frac{H\nu\cdot(x-y)}{r^n}g\,dA
-\frac1{n\wn}\int_\mathcal S\frac{(x-y)\cdot\delta g}{r^n}\,dA,\\
g(y)&\le\frac1{4\wn}\int_\mathcal SgH\,r^{2-n}\,dA
-\frac1{n\wn}\int_\mathcal S\frac{(x-y)\cdot\delta g}{r^n}\,dA.
\end{aligned}
\tag{16.40}
\end{equation}
The final integral may be rewritten using (16.13).
\end{remark}

\section*{16.2. Interior Gradient Bounds}

Let $\mathcal G$ be the graph of $u\in C^2(\Omega)$, put
$v=\sqrt{1+|Du|^2}$, and let $\nu=v^{-1}(-Du,1)$ be its upper normal.
Then

\[
\tag{16.41}
H=-\frac1nD_i\nu_i
=\frac1nD_i\!\left(\frac{D_iu}{v}\right).
\]
Equivalently,
\[
\tag{16.42}
\int_\Omega(\nu\cdot D\varphi-nH\varphi)\,dx'=0
\]
for $\varphi\in C_0^1(\Omega)$.

\[
\tag{16.43}
\int_\Omega(D_k\nu\cdot D\varphi-nD_kH\varphi)\,dx'=0.
\]
Here $H\in C^1(\Omega)$.  Testing (16.43) by $\nu_k\varphi$ gives
\[
\tag{16.44}
\int_\Omega D_k\nu_iD_i\nu_k\varphi\,dx'
+\int_\Omega(\nu_kD_k\nu_iD_i\varphi
-n\nu_kD_kH\varphi)\,dx'=0.
\]
\[
\nu_kD_k\nu_i=-\frac{\nu_k}{v}(\delta_{ij}-\nu_i\nu_j)D_{jk}u
=-v(\delta_{ij}-\nu_i\nu_j)D_j(v^{-1})
=-v\delta_i\nu_{n+1}.
\]
Consequently

\[
\tag{16.45}
\int_\Omega\mathcal C^2\varphi\,dx'
-\int_\Omega v(\delta_i\nu_{n+1}D_i\varphi
-n\delta_{n+1}H\varphi)\,dx'=0,
\]
where
\[
\mathcal C^2=\sum_{i,j=1}^nD_i\nu_jD_j\nu_i
=\sum_{i,j=1}^{n+1}(\delta_i\nu_j)^2.
\]
Lifting test functions to $\mathcal U=\Omega\times\mathbb R$ yields
\[
\tag{16.46}
\int_\mathcal G(\mathcal C^2\nu_{n+1}\varphi
-\delta_i\nu_{n+1}\delta_i\varphi
+n\delta_{n+1}H\varphi)\,dA=0.
\]
Set
\[
\tag{16.47}
w=\log v=-\log\nu_{n+1}.
\]
Replacing $\varphi$ by $v\varphi$ in (16.46) gives

\[
\tag{16.48}
\int_\mathcal G(\delta w\cdot\delta\varphi+|\delta w|^2\varphi
-n\nu\cdot DH\varphi)\,dA\le0
\]
for nonnegative $\varphi\in C_0^1(\mathcal U)$, or equivalently
\[
\tag{16.49}
\Delta w\ge|\delta w|^2-n\nu\cdot DH.
\]
For $H=0$, $w$ is weakly subharmonic.  In general, Lemma 16.3 gives

\begin{align*}
w(y') &\leq \frac{1}{\omega_n R^n} \int_{\mathcal{G}_{R(y)}} w\, dA + \frac{1}{4\omega_n} \int_{\mathcal{G}_{R(y)}} w H^2 r^{2}(r^{-n} - R^{-n})\, dA \tag{16.50}\\
&\qquad + \frac{1}{\omega_n} \int_{\mathcal{G}_{R(y)}} \psi(r)\, \nu \cdot DH\, dA
\end{align*}

provided $R < \operatorname{dist}(y', \partial\Omega)$, $y = (y', u(y'))$ and $\psi$ is given by (16.27). Setting

\begin{equation}\tag{16.51}
H_0 = \sup_{\Omega} |H|,\qquad
H_1 = \sup_{\Omega} (\nu \cdot DH)^+ \leq \sup_{\Omega} |DH|,
\end{equation}

we have from (16.50)

\begin{align*}
w(y') &\leq \frac{1}{\omega_n R^n} \int_{\mathcal{G}_R} w\, dA + \frac{H_0^2}{4\omega_n} \int_{\mathcal{G}_R} w r^2(r^{-n} - R^{-n})\, dA \\
&\qquad + \frac{H_1}{\omega_n} \int_{\mathcal{G}_R} \psi(r)\, dA \\
&= \frac{1}{\omega_n R^n} \int_{\mathcal{G}_R} w\, dA + \frac{n H_0^2}{4\omega_n} \int_0^R \rho^{-n-1} \int_{\mathcal{G}_\rho} w r^2\, dA\, d\rho \\
&\qquad + \frac{H_1}{\omega_n} \int_0^R \rho(\rho^{-n} - R^{-n}) A(\mathcal{G}_\rho)\, d\rho \qquad \text{(by Fubini's theorem)} \tag{16.52}\\
&\leq \frac{1}{\omega_n R^n} \int_{\mathcal{G}_R} w\, dA + \frac{n H_0^2}{4\omega_n} \int_0^R \rho^{1-n} \int_{\mathcal{G}_\rho} w\, dA\, d\rho \\
&\qquad + \frac{H_1}{\omega_n} \int_0^R \rho^{1-n} A(\mathcal{G}_\rho)\, d\rho.
\end{align*}

It remains to bound

\[
A(\mathcal{G}_\rho) \text{ and } \int_{\mathcal{G}_\rho} w\, dA \text{ for } 0 < \rho \leq R.
\]



For the area term, assume $3R<\dist(y',\partial\Omega)$ and normalize
$y'=0$, $u(y')=0$. For $\rho\le R$, define

\[
u_\rho =
\begin{cases}
\rho & \text{for } u \geq \rho \\
u & \text{for } -\rho \leq u \leq \rho \\
-\rho & \text{for } u \leq -\rho
\end{cases}
\]

Use the test function

\[
\varphi = \eta u_\rho
\]

in (16.42), where $\eta$ is uniformly Lipschitz, equals one on
$|x'|<\rho$, vanishes on $|x'|>2\rho$, and satisfies $|D\eta|\le1/\rho$.
The identity extends to such compactly supported $W_0^{1,1}$ tests and gives

\[
\int_{|x'|,\,|u|<\rho} \frac{|Du|^2}{v}\, dx' \leq \rho \int_{|x'|<2\rho} \left(|D\eta| + n\eta |H|\right) dx'.
\]

Consequently

\begin{align*}
A(\mathcal{G}_\rho) &= \int_{|x'|^2 + u^2 < \rho^2} v\, dx' \\
&\leq \int_{|x'|,\,|u|<\rho} v\, dx' \tag{16.53} \\
&\leq C(n)\left\{\rho^n + \rho \int_{|x'|<2\rho} |H|\eta\, dx'\right\} \\
&\leq C(n)\rho^n(1 + H_0 R).
\end{align*}

For the logarithmic term, substitute

\[
\varphi = \eta w(u_\rho + \rho)
\]







in (16.42), with $\eta$ as above. Then

\[
\int_{|x'|,\,|u|<\rho} \frac{w|Du|^2}{v}\,\dxp
\le 2\rho \int_{|x'|<2\rho,\,u>-\rho} (w|D\eta|+\eta|Dw|+n\eta|H|w)\,\dxp.
\]

Replacing $\varphi$ by $\varphi^2$ in (16.48) gives

\[
\int_{\mathcal{G}} \varphi^2|\delta w|^2\,dA
\le -2\int_{\mathcal{G}} \varphi\,\delta w\cdot\delta\varphi\,dA
- n\int_{\mathcal{G}} \varphi^2\,\nu\cdot DH\,dA
\]

for $\varphi\in C_0^1(\Omega\times\mathbb R)$. Cauchy's inequality yields

\[
\int_{\mathcal{G}} \varphi^2|\delta w|^2\,dA
\le 4\int_{\mathcal{G}} |\delta\varphi|^2\,dA + nH_1 \int_{\mathcal{G}} \varphi^2\,dA.
\]

Choose

\[
\varphi=\eta\tau(x_{n+1}),
\]

where $0\le \tau\le 1$, $\tau\equiv 1$ in $\bigl(-\rho,\sup_{|x'|<2\rho}u\bigr)$, $\tau\equiv 0$ outside $\bigl(-2\rho,\rho+\sup_{|x'|<2\rho}u\bigr)$, $|\tau'|<2/\rho$
and $\eta$ is as above. We then have

\[
\int_{\mathcal{G}} \varphi^2|\delta w|^2\,dA \le (8\rho^{-2}+nH_1)\,A(\mathcal{G}\cap\supp\varphi).
\]

Equation (16.8) then gives

\[
\int_{|x'|<2\rho,\,u>-\rho}\eta|Dw|\,\dxp
\le \int_{\mathcal{G}} \varphi|\delta w|\,dA
\]

\[
\le (8\rho^{-2}+nH_1)^{1/2}A(\mathcal{G}\cap\supp\varphi)
\quad\text{by Schwarz's inequality.}
\]

\[
\le (8+nH_1R^2)^{1/2}\rho^{-1}
\int_{|x'|<2\rho,\,u>-\rho} v\,\dxp.
\]

Since $w\le v$, we also have

\[
\int_{|x'|<2\rho,\,u>-\rho} w\,\dxp \le \int_{|x'|<2\rho,\,u>-\rho} v\,\dxp.
\]

Finally, use

\[
\varphi=\eta\max\{u+\rho,0\}
\]







in (16.42), where $\eta=1$ on $|x'|<2\rho$, $\eta=0$ on
$|x'|>3\rho$, and $|D\eta|<1/\rho$. This gives

\[
\int_{|x'|<2\rho,\,u>-\rho} v\,\dxp \le C(n)\rho^n(1+H_0R)\bigl(1+\rho^{-1}\sup_{|x'|<3\rho}u\bigr).
\]

Combining these estimates gives

\[
\int_{\mathcal{G}_\rho} w\,\dA \le \int_{|x'|,\,|u|<\rho} wv\,\dxp
\]

\begin{equation}\tag{16.54}
\le C(n)\rho^n(1+H_0R)(1+H_1R^2)^{1/2}\bigl(1+\rho^{-1}\sup_{\Omega}u\bigr).
\end{equation}

\begin{theorem}[Interior gradient bound]\label{gilbarg-trudinger.thm.16.5.interior-gradient-bound}\dcref{ch16:16.5}\group{chapter-16}\uses{gilbarg-trudinger.lem.16.3.potential-estimate}
Let $\Omega\subset\mathbb R^n$ be a domain and $u\in C^2(\Omega)$. For $y'\in\Omega$,

\begin{equation}\tag{16.55}
|Du(y')|\le C_1 \exp\{C_2 \sup_{\Omega}(u-u(y'))/d\},
\end{equation}

where $d=\dist(y',\partial\Omega)$ and $C_i=C_i(n,dH_0,d^2H_1)$, with $H_0,H_1$ as in (16.51).
\end{theorem}
\begin{proof}
\sketch Apply Lemma 16.3 to $w=\log\sqrt{1+|Du|^2}$ using (16.48)--(16.50). The truncation tests in (16.42) give the area bound (16.53) and logarithmic integral bound (16.54); substitution into (16.52), followed by exponentiation, gives (16.55).
\end{proof}

\begin{corollary}[Gradient bound for nonnegative solutions]\label{gilbarg-trudinger.cor.16.6.nonnegative-gradient-bound}\dcref{ch16:16.6}\group{chapter-16}\uses{gilbarg-trudinger.thm.16.5.interior-gradient-bound}
Under the hypotheses of Theorem 16.5, if $u\geq0$, then for every $y'\in\Omega$,

\begin{equation}\tag{16.56}
|Du(y')|\le C_1 \exp\{C_2u(y')/d\}
\end{equation}

with $C_1,C_2,d$ as in Theorem 16.5.
\end{corollary}

\begin{corollary}[Higher derivative interior estimate]\label{gilbarg-trudinger.cor.16.7.higher-derivative-estimate}\dcref{ch16:16.7}\group{chapter-16}\uses{gilbarg-trudinger.thm.16.5.interior-gradient-bound}
Let $u\in C^2(\Omega)$ and suppose its graph has mean curvature $H\in C^k(\Omega)$ for $k\geq1$. Then $u\in C^{k+1}(\Omega)$, and for $y\in\Omega$ and $|\beta|=k+1$,

\begin{equation}
|D^\beta u(y)|\le C.
\tag{16.57}
\end{equation}

where $C=C(n,k,|H|_{k,\Omega},d,\sup_\Omega|u|)$ and $d=\dist(y,\partial\Omega)$.
\end{corollary}

\section*{16.3. Application to the Dirichlet Problem}

\begin{theorem}[Continuous Dirichlet data for the minimal-surface equation]\label{gilbarg-trudinger.thm.16.8.minimal-surface-dirichlet}\dcref{ch16:16.8}\group{chapter-16}\uses{gilbarg-trudinger.cor.16.7.higher-derivative-estimate}
Let $\Omega\subset\mathbb R^n$ be a bounded $C^2$ domain. The problem
\[
\mathcal Mu=0\quad\text{in }\Omega,\qquad u=\varphi\quad\text{on }\partial\Omega
\]
is solvable for every $\varphi\in C^0(\partial\Omega)$ if and only if the mean curvature of $\partial\Omega$ is nonnegative everywhere.
\end{theorem}
\begin{proof}
\sketch First assume $\partial\Omega\in C^{2,\alpha}$ and approximate $\varphi$ uniformly by $C^{2,\alpha}$ boundary data $\varphi_m$. The smooth Dirichlet theorem and comparison principle give solutions $u_m$ satisfying

\begin{equation*}
\sup_{\Omega} |u_{m_1} - u_{m_2}| \leq \sup_{\partial\Omega} |\varphi_{m_1} - \varphi_{m_2}| \to 0 \qquad \text{as } m_1, m_2 \to \infty.
\end{equation*}

Thus $u_m$ converges uniformly; Corollary 16.7 and compactness give a $C^2$ interior solution with the required trace. Approximation of the domain handles $C^2$ boundaries, while the boundary barrier obstruction gives necessity.
\end{proof}

\begin{theorem}[Conditional prescribed-mean-curvature solvability]\label{gilbarg-trudinger.thm.16.9.conditional-dirichlet-solvability}\dcref{ch16:16.9}\group{chapter-16}
Let $\Omega\subset\mathbb R^n$ be bounded with $\partial\Omega\in C^{2,\alpha}$, $0<\alpha<1$, let $\varphi\in C^{2,\alpha}(\overline\Omega)$, and let $H\in C^1(\overline\Omega)$. Assume the boundary mean curvature $H'$ satisfies

\begin{equation}
H'(y)\ge\frac n{n-1}|H(y)|.
\tag{16.58}
\end{equation}
at each $y\in\partial\Omega$. Then there is a unique
$u\in C^{2,\alpha}(\overline\Omega)$ satisfying
equation (16.1) in $\Omega$ and $u=\varphi$ on $\partial\Omega$ provided the family of solutions of the Dirichlet
problems

\begin{equation}\tag{16.59}
\M u = \sigma nH(x)\bigl(1+|Du|^2\bigr)^{3/2}\ \text{in }\Omega,\qquad u=\sigma\varphi\ \text{on }\partial\Omega,
\end{equation}

is uniformly bounded in $\Omega$.
\end{theorem}

For a solution, (16.42) implies

\[
\left|\int_\Omega H\eta\,dx\right|\leq \frac{1}{n}\int_\Omega |\nu\cdot D\eta|\,dx
\]
\[
\leq \frac{1}{n}\sup_\Omega \frac{|Du|}{\sqrt{1+|Du|^2}}\int_\Omega |D\eta|\,dx
\]

Writing

\[
1-\varepsilon_0=\sup_\Omega \frac{|Du|}{\sqrt{1+|Du|^2}},
\]

gives the necessary condition
\begin{equation}\tag{16.60}
\left|\int_\Omega H\eta\,dx\right|\leq \frac{1-\varepsilon_0}{n}\int_\Omega |D\eta|\,dx,
\end{equation}
for every $\eta\in C_0^1(\Omega)$ and some $\varepsilon_0>0$.
Conversely, (16.60) supplies the height estimate required in Theorem 16.9.

\begin{theorem}[Sharp prescribed-mean-curvature Dirichlet theorem]\label{gilbarg-trudinger.thm.16.10.prescribed-mean-curvature-dirichlet}\dcref{ch16:16.10}\group{chapter-16}\uses{gilbarg-trudinger.thm.16.9.conditional-dirichlet-solvability}
Let $\Omega$ be a bounded domain in $\mathbb{R}^n$ with $\partial\Omega \in C^{2,\alpha}$ for some $\alpha$,
$0<\alpha<1$, and let $\varphi \in C^{2,\alpha}(\overline{\Omega})$. Let $H \in C^1(\overline{\Omega})$ satisfy (16.58) and (16.60). Then
the Dirichlet problem $\M u = nH(1+|Du|^2)^{3/2},\ u=\varphi\ \text{on }\partial\Omega,$ is uniquely solvable
for $u \in C^{2,\alpha}(\overline{\Omega})$. Furthermore, if we only assume that $\varphi \in C^0(\partial\Omega)$, the problem is
uniquely solvable for $u \in C^0(\overline{\Omega}) \cap C^2(\Omega)$.
\end{theorem}
\begin{proof}
\sketch Condition (16.60) supplies the uniform height estimate required in Theorem 16.9. Smooth data give the classical solution; uniform approximation and Theorem 16.8's compactness argument give the continuous-boundary solution. Comparison gives uniqueness.
\end{proof}

Condition (16.60) follows in particular from

\begin{equation}\tag{16.61}
\int_\Omega |H|^n\,dx < \omega_n,
\end{equation}
For a solution only in $C^0(\overline\Omega)\cap C^2(\Omega)$, condition
(16.60) may be used with $\varepsilon_0=0$.

For constant $H$, let $\Omega_1$ be the set on which the nearest boundary
point is unique.  Condition (16.58) and the distance-function identities give

\[
\Delta d \le -n|H|
\]

in $\Omega_1$, where $d(x)=\operatorname{dist}(x,\partial\Omega)$.  Set

\[
 v(x)=\sup_{\partial\Omega}|u|+\frac{e^{\mu\delta}}{\mu}\bigl(1-e^{-\mu d(x)}\bigr),
\]

where $\delta=\operatorname{diam}\Omega$ and $\mu=1+n|H|$.  Then

\[
\Mv v = \bigl[-\mu + (1+|Dv|^2)\,\Delta d\bigr]e^{\mu(\delta-d)}
\]
\[
\le -\bigl[\mu+n|H|(1+|Dv|^2)\bigr]e^{\mu(\delta-d)}
\]
\[
\le -n|H|(1+|Dv|^2)^{3/2},
\]

Thus $v$ is a supersolution on $\Omega_1$.  Comparison, followed along a
normal segment at points outside $\Omega_1$, gives

\begin{equation}
\sup_\Omega|u|\le\sup_{\partial\Omega}|u|+(e^{n\delta}-1)/\mu.
\tag{16.62}
\end{equation}

\begin{theorem}[Constant-mean-curvature Dirichlet criterion]\label{gilbarg-trudinger.thm.16.11.constant-mean-curvature-dirichlet}\dcref{ch16:16.11}\group{chapter-16}\uses{gilbarg-trudinger.thm.16.9.conditional-dirichlet-solvability}
Let $\Omega$ be a bounded $C^2$ domain in $\mathbb{R}^n$. For constant $H$ and arbitrary $\varphi\in C^0(\partial\Omega)$, the problem
\[
\Mv u=nH(1+|Du|^2)^{3/2},\qquad u=\varphi \text{ on } \partial\Omega,
\]
is solvable if and only if the boundary mean curvature satisfies $H'\ge n|H|/(n-1)$ everywhere.
\end{theorem}
\begin{proof}
\sketch For $d(x)=\dist(x,\partial\Omega)$ on the unique-nearest-point region, use $\Delta d\leq-n|H|$ and the exponential barrier displayed above to obtain (16.62). Theorem 16.9 then yields existence; the boundary obstruction gives necessity and comparison gives uniqueness.
\end{proof}

\section*{16.4. Equations in Two Independent Variables}

For $\Omega\subset\mathbb R^2$, consider

\begin{equation}\tag{16.63}
Qu=a^{ij}(x,u,Du)D_{ij}u+b(x,u,Du)=0,
\end{equation}

and assume, for fixed $\gamma\ge1$ and $\mu\ge0$, that

\begin{equation}\tag{16.64}
|\xi|^2-\frac{(p\cdot\xi)^2}{1+|p|^2}\le a^{ij}(x,z,p)\xi_i\xi_j\le \gamma\left[|\xi|^2-\frac{(p\cdot\xi)^2}{1+|p|^2}\right]
\end{equation}

for every $(x,z,p)$ and $\xi\in\mathbb R^2$, and

\begin{equation}\tag{16.65}
|b(x,z,p)|\le \mu\sqrt{1+|p|^2}
\end{equation}

for every $(x,z,p)$.  The minimal-surface equation has

\begin{equation*}
a^{ij}(x,z,p)=\delta^{ij}-\frac{p_ip_j}{1+|p|^2},\qquad b=0;
\end{equation*}

so $\gamma=1$ and $\mu=0$.

\begin{proposition}[Geometric characterization of mean-curvature-type equations]
\label{gilbarg-trudinger.prop.16.mean-curvature-type-characterization}\dcref{ch16:16.63'--16.65'}\group{chapter-16}
Let $u\in C^2(\Omega)$ have graph

\[
\mathcal{G}=\{(x,z)\in\mathbb{R}^3\mid x\in\Omega,\ z=u(x)\}.
\]

There are coefficients satisfying (16.63)--(16.65) if and only if the
principal curvatures satisfy

\begin{equation}\tag{16.63'}
\alpha_1\kappa_1+\alpha_2\kappa_2+\beta=0,
\end{equation}

where

\begin{equation}\tag{16.64'}
1\le \alpha_i\le \gamma,\qquad i=1,2,
\end{equation}

\begin{equation}\tag{16.65'}
|\beta|\le \mu.
\end{equation}
\end{proposition}
\begin{proof}
\sketch Let $d(x,z)$ be the signed distance to $\mathcal G$, positive when
$z>u(x)$.  On the graph, where $X=(x,u(x))$, the chain rule gives

\[
D_i d(X)+D_i u(x)D_3 d(X)=0
\]

and

\begin{equation}
D_{ij}d(X)+D_i u(x)D_3 d(X)+D_j u(x)D_{3i}d(X)+D_i u(x)D_j u(x)D_{33}d(X)
+D_3 d(X)D_{ij}u(x)=0,
\tag{16.66}
\end{equation}

$i,j=1,2$, where $X=(x,u(x))$. Since $D_3 d(X)=v^{-1}$, $v=\sqrt{1+|Du(x)|^2}$, (16.63)
then implies

\begin{equation}
\sum_{i,j=1}^3 a_*^{ij}(x)D_{ij}d(X)+v^{-1}b_*(x)=0,
\tag{16.67}
\end{equation}

where $b_*(x)=b(x,u(x),Du(x))$ and where the $3\times 3$ matrix $[a_*^{ij}(x)]$ is defined by
setting $a_*^{ij}(x)=a^{ij}(x,u(x),Du(x))$ for $i,j=1,2$, and

\[
a_*^{i3}(x)=a_*^{3i}(x)=\sum_{j=1}^2 D_j u(x)a_*^{ij}(x),\qquad i=1,2,
\]
\[
a_*^{33}(x)=\sum_{i,j=1}^2 D_i u(x)D_j u(x)a_*^{ij}(x).
\]

Equivalently,

\[
\sum_{j=1}^3 a_*^{ji}(x)\nu_j=\sum_{j=1}^3 a_*^{ij}(x)\nu_j=0,\qquad i=1,2,3,
\]

where $\nu=v^{-1}(-Du(x),1)=Dd(X)$ is the upward unit normal of $\Ssurface$.
Let $q$ be an isometry to principal coordinates at $X$, so that
$q^t[D_{ij}d(X)]q=\operatorname{diag}[\kappa_1,\kappa_2,0]$,
where $\kappa_1,\kappa_2$ are the principal curvatures. Thus (16.67) has the
form (16.63'), with $\alpha_1,\alpha_2$ the first two diagonal entries of
$q^t[a_*^{ij}(x)]q$ and $\beta=v^{-1}b_*(x)$. Condition (16.65') follows
from (16.65), while (16.64') follows from

\[
\sum_{i,j=1}^3 a_*^{ij}(x)\xi_i\xi_j=\sum_{i,j=1}^2 a_*^{ij}(x)(\xi_i+\xi_3D_i u(x))(\xi_j+\xi_3D_j u(x)),\qquad \xi\in \R^3,
\]



and (16.64):
\[
|\xi'|^2 \le \sum_{i,j=1}^3 a_*^{ij}(x)\xi_i\xi_j \le \gamma|\xi'|^2, \qquad \xi' = \xi - (\nu\cdot \xi)\nu,\ \nu = v^{-1}(-Du(x),1).
\]
This proves (16.64').

Conversely, assume (16.63')--(16.65') at $X=(x,u(x))\in\G$ and set
$[a_*^{ij}(x)]=q\,\mathrm{diag}[\alpha_1,\alpha_2,0]q$ and
$b_*(x)=v\beta$. Then (16.67) holds; using
\[
\sum_{j=1}^3 a_*^{ij}(x)\nu_j=0, \qquad i=1,2,3,
\]
an application of (16.66) yields
\[
\sum_{i,j=1}^2 a_*^{ij}(x)D_{ij}u+b_*(x)=0.
\]

Define, for $i,j=1,2$,
\[
a^{ij}(x,z,p)=\begin{cases}
 a_*^{ij}(x) & \text{if } z=u(x)\text{ and }p=Du(x)\\
 \delta_{ij}-\dfrac{p_ip_j}{1+|p|^2} & \text{otherwise,}
\end{cases}
\]
and
\[
b(x,z,p)=\begin{cases}
 b_*(x) & \text{if } z=u(x)\text{ and }p=Du(x)\\
 0 & \text{otherwise.}
\end{cases}
\]
\end{proof}

\begin{theorem}[Bernstein theorem for mean-curvature-type equations]\label{gilbarg-trudinger.thm.16.12.bernstein}\dcref{ch16:16.12}\group{chapter-16}\uses{gilbarg-trudinger.cor.16.19.quasiconformal-gauss-map-bernstein}
Let $u\in C^2(\mathbb R^2)$ solve (16.63) on $\mathbb R^2$, with $b\equiv0$, and assume (16.64). Then $u$ is affine.
\end{theorem}
\begin{proof}
\sketch By (16.89), the Gauss map of the graph is $(K,0)$-quasiconformal with $K=-\gamma$. Corollary 16.19 makes the normal constant, hence $Du$ is constant.
\end{proof}



\section*{16.5. Quasiconformal Mappings}

Let $\Sc,\Tc\subset\mathbb R^3$ be oriented $C^2$ hypersurfaces with
unit normals $\nu,\mu$.  Fix $Y\in\Sc$, write
$\Sc_\rho(Y)=\Sc\cap B_\rho(Y)$, and take
$\Sc_R(Y)\Subset\Sc$.  For a $C^1$ vector field $V$ and a regular
$\Gc\Subset\Sc$, Stokes' theorem is

\begin{equation}\tag{16.68}
\int_\Gc\nu\cdot(D\times V)\,dA
=\int_{\partial\Gc}V_i\,dx_i
=\int_{\partial\Gc}\mathbf t\cdot V\,ds.
\end{equation}
Taking $V=fDg$ gives, for $f,g\in C^1(\Sc)$,
\begin{equation}\tag{16.69}
\int_\Gc\nu\cdot(\delta f\times\delta g)\,dA
=\int_{\partial\Gc}f\,dg
=\int_{\partial\Gc}f\frac{dg}{ds}\,ds.
\end{equation}
Assume a $C^1$ vector field $\omega$ near $\Tc$ satisfies

\begin{equation}\tag{16.70}
\sup_{\Tc}|\omega|+\sup_{\Tc}|D\omega|\leq \Lambda_0\qquad \text{and}\qquad \mu\cdot D\times \omega \equiv 1\text{ on }\Tc.
\end{equation}

Then

\begin{equation}\tag{16.71}
A(\Gc)=\int_{\partial\Gc}\omega_i\,dx_i,
\end{equation}

For the upper unit hemisphere, take $\mu(X)=X$ and

\[
\omega(X)=(-x_2/(1+x_3),\ x_1/(1+x_3),\ 0).
\]

For a $C^1$ map $\varphi:\Gc\to\Tc$, define its signed area
magnification by

\begin{equation}\tag{16.72}
J(\varphi)=\nu\cdot
\delta(\omega_i\circ\varphi)\times\delta\varphi_i.
\end{equation}
Indeed, if $\varphi$ is one-to-one on a regular region $\Ec\subset\Gc$,
(16.69)--(16.71) give

\[
A(\varphi(\Ec))
=\pm\int_\Ec J(\varphi)\,dA,
\]
with sign determined by orientation.  The definition is invariant under
orientation-preserving isometries: if $P,Q$ are such isometries and

\[
\det P=\det Q=1,
\]

and if we define

\[
\widetilde{\Sc}=\{P(X-Y)\mid X\in\Sc\},\ \widetilde{\Tc}=\{Q(X-\varphi(Y))\mid X\in\Tc\},
\]
\[
\widetilde{\nu}=P\circ\nu\circ P^{-1},\quad
\widetilde{\mu}=Q\circ\mu\circ Q^{-1},\quad
\widetilde{\omega}=Q\circ\omega\circ Q^{-1},\quad
\widetilde{\varphi}=Q\circ\varphi\circ P^{-1},
\]

then

\begin{equation}\tag{16.73}
\nu\cdot(\widetilde\delta\omega_i\circ\varphi)
\times(\delta\varphi_i)(Y)
=\widetilde\nu\cdot
(\widetilde\delta\widetilde\omega_i\circ\widetilde\varphi)
\times(\widetilde\delta\widetilde\varphi_i)(0),
\end{equation}

and $\widetilde\delta$ is the tangential gradient on $\widetilde\Sc$.
Orthogonal invariance of the scalar triple product proves (16.73), while
(16.70) gives

\[
\widetilde{\mu}\cdot D\times \widetilde{\omega}\equiv 1\quad \text{on }\widetilde{\Sc}.
\]

If the isometries $P, Q$ above are chosen so that

\[
\widetilde\nu(0)=\widetilde\mu(0)=(0,0,1),
\]

and if we introduce new coordinates

\[
(\zeta,\zeta_3)=(\zeta_1,\zeta_2,\zeta_3)=P(X-Y),
\]

then $\widetilde{\Tc}$ can be represented near $0$ in the form

\[
\zeta_3=\widetilde{u}(\zeta),\ \zeta\in\Uc,
\]

with $D\widetilde u(0)=0$.  Tangency gives

\[
\widetilde{\delta}_j\widetilde{\varphi}_3(0)=(0,0,1)\cdot\widetilde{\delta}_j\widetilde{\varphi}(0)
=\widetilde{\mu}(0)\cdot\widetilde{\delta}_j\widetilde{\varphi}(0)=0,\qquad j=1,2,3
\]

and hence

\[
\widetilde{\delta}\widetilde{\varphi}_3(0)=0.
\]



Thus if
\[
\psi : \mathcal U \to \mathbb R^2
\]
is defined by
\[
\psi(\zeta) = [\tilde{\varphi}_1(\zeta,\tilde{u}(\zeta)), \tilde{\varphi}_2(\zeta,\tilde{u}(\zeta))], \qquad \zeta \in \mathcal U,
\]
then $\psi$ approximates $\tilde{\varphi}$ near $0$ in the sense that
\[
\bigl|(\psi(\zeta),0) - \tilde{\varphi}(\zeta,\tilde{u}(\zeta))\bigr| = o(|\zeta|) \text{ as } |\zeta| \to 0, \qquad \zeta \in \mathcal U.
\]

Equations (16.73) and the tangential projection formula give

\begin{equation}
J(\varphi)(Y) = D_1\psi_1(0)D_2\psi_2(0) - D_1\psi_2(0)D_2\psi_1(0);
\tag{16.74}
\end{equation}

and

\begin{equation}
|\delta\varphi(Y)|^2 = |\tilde{\delta}\tilde{\varphi}(0)|^2 = |D\psi(0)|^2.
\tag{16.75}
\end{equation}

A map $\varphi:\mathcal G\to\mathcal Z$ is $(K,K')$-quasiconformal when

\begin{equation}
|\delta\varphi(X)|^2 \leq 2KJ(\varphi)(X) + K'.
\tag{16.76}
\end{equation}

at every $X\in\mathcal G$, where $K\in\mathbb R$ and $K'\ge0$.
If $K'=0$ and $\delta\varphi\not\equiv0$, then $|K|\ge1$.

\begin{lemma}[Surface coarea identity]
\label{gilbarg-trudinger.lem.16.surface-coarea}\dcref{ch16:16.77--16.78}\group{chapter-16}
For $g\in C^1(\mathcal G_\rho)$ and almost every $\rho\in(0,R)$,

\begin{equation}
\int_{\partial \mathcal G_\rho} g\, ds = D_\rho \int_{\mathcal G_\rho} g|\delta r|\, dA
\tag{16.77}
\end{equation}

where
\[
r(X) = |X - Y|, \qquad X \in \mathbb R^3.
\]
For such regular radii,

\begin{equation}\tag{16.78}
\partial \widetilde{\mathscr{S}}_{\rho}
= \widetilde{\mathscr{S}}_{\rho} \cap \{X \in \mathbb{R}^{3}\mid |X-Y|=\rho\}
= \bigcup_{j=1}^{n(\rho)} \Gamma_{\rho}^{(j)}
\end{equation}

where the $\Gamma_\rho^{(j)}$ are simple closed $C^2$ curves.
\end{lemma}
\begin{proof}
\sketch Sard's theorem gives $\delta r\ne0$ on the boundary for almost every
$\rho$.  On the annulus
$\widetilde{\mathscr G}=\widetilde{\mathscr S}_\rho
\setminus\widetilde{\mathscr S}_{\rho-\varepsilon}$, orient the tangent by
$\mathbf t=\nu\times\delta r/|\delta r|$ on the outer boundary and with
the opposite sign on the inner boundary.  Extend
$\nu\times\delta r/|\delta r|$ to a $C^1$ field $\mathbf F$ and apply
(16.68) with the cutoff field
$\varepsilon^{-1}(r-\rho+\varepsilon)g\mathbf F$.
The boundary term is the difference quotient of
$\int_{\mathcal G_\rho}g|\delta r|\,dA$; all remaining annular terms tend
to zero.  Letting $\varepsilon\downarrow0$ gives (16.77).
\end{proof}
The main H\"older estimate will be a consequence of estimates for the integral
\[
\D(\rho; Z),
\]
which is defined for $Z \in \Gscr$ and $\rho > 0$ by
\[
\D(\rho; Z)=\int_{\Sscr_{\rho}(Z)} |\delta \varphi|^2\, dA.
\]

(cf. the Dirichlet integral used in Chapter 12).
$\Lambda_1$ denotes a constant such that

\begin{equation}\tag{16.79}
\A(\Sscr_{3R/4}(Y)) \le \Lambda_1 (3R/4)^2.
\end{equation}

\begin{lemma}[Dirichlet integral bound for quasiconformal maps]\label{gilbarg-trudinger.lem.16.13.dirichlet-bound}\dcref{ch16:16.13}\group{chapter-16}
Suppose $\varphi:\Gscr\to\Tscr$ is $(K,K')$-quasiconformal and $C^1$. Then

\begin{equation}\tag{16.80}
\D(R/2; Y) \le C
\end{equation}

where $C=C(\Lambda_0,K,K'R^2,\Lambda_1)$.
\end{lemma}

\begin{proof}
\sketch For regular radii, Stokes' theorem gives
\begin{equation}
\int_{\mathcal S_\rho}J(\varphi)\,dA
=\sum_{j=1}^{n(\rho)}\int_{\Gamma_\rho^{(j)}}
(\omega_i\circ\varphi)\frac{d\varphi_i}{ds}\,ds.
\tag{16.81}
\end{equation}
Together with (16.70), (16.76), and coarea (16.77), this bounds the
Dirichlet integral by a product of radial derivatives.  Introduce the
augmented functions $g=f+\rho R$ and
$\mathcal E=\mathcal D+\rho/R$; the resulting inequality is
\begin{equation}
-\frac d{d\rho}\left(\frac1{\mathcal E(\rho)}\right)
\ge\frac1{Cg'(\rho)}.
\tag{16.82}
\end{equation}
H\"older's inequality on $(R/2,3R/4)$ and (16.79) then give (16.80).
\end{proof}

\begin{theorem}[Morrey decay of the Dirichlet integral]\label{gilbarg-trudinger.thm.16.14.dirichlet-decay}\dcref{ch16:16.14}\group{chapter-16}\uses{gilbarg-trudinger.lem.16.13.dirichlet-bound}
Suppose $\varphi:\Sscr\to\Tscr$ is $(K,K')$-quasiconformal and $C^1$, and assume
\begin{equation}\tag{16.83}
\int_{\Sscr_{R/2}(Y)}H^2\,dA\leq\Lambda_2.
\end{equation}
Then
\begin{equation}\tag{16.84}
\Dop(\rho;Z)\leq C(\rho/R)^\alpha
\end{equation}
for $Z\in\mathcal G_{R/4}(Y)$ and $0<\rho<R/4$, where $C,\alpha>0$ depend only on $\Lambda_0,K,K'R^2,\Lambda_1,\Lambda_2$.
\end{theorem}
\begin{proof}
\sketch Apply (16.81) on each boundary component after subtracting a base value of $\omega\circ\varphi$. Cauchy--Schwarz, (16.70), (16.76), and (16.77) imply $\mathcal D(\rho)\leq C\rho\mathcal D'(\rho)+C'K'\rho^2$. Adding $(\rho/R)^2$ and integrating the resulting logarithmic inequality gives (16.84).
\end{proof}


\begin{theorem}[H\"older continuity of quasiconformal maps]\label{gilbarg-trudinger.thm.16.15.quasiconformal-holder}\dcref{ch16:16.15}\group{chapter-16}\uses{gilbarg-trudinger.thm.16.14.dirichlet-decay}\uses{gilbarg-trudinger.lem.16.4.surface-morrey-estimate}
Suppose $\varphi:\G\to\Sigma$ is $(K,K')$-quasiconformal and $C^1$. Then

\begin{equation}\tag{16.85}
\sup_{X\in \Sset_\rho^\ast(Y)} |\varphi(X)-\varphi(Y)|\leq C(\rho/R)^\alpha,\qquad \rho\in(0,R/4),
\end{equation}






where $C,\alpha>0$ depend only on $\Lambda_0,K,K'R^2,\Lambda_1,\Lambda_2$, and $\Sg_\rho^*(Y)$ is the component containing $Y$.
\end{theorem}
\begin{proof}
\sketch H\"older's inequality, (16.84), and the area estimate (16.34) give the hypothesis of Lemma 16.4 for each component, with exponent $\alpha/2$; (16.38) yields (16.85).
\end{proof}

\section*{16.6.\ Graphs with Quasiconformal Gauss Map}

Let $\Sg$ be the graph of $u\in C^2(\Omega)$, fix
$Y=(y,u(y))\in\Sg$, and suppose $B_R(y)\subset\Omega$.  Its upward unit
normal is

\[
\unitv(x,z) \equiv \unitv(x) = \left(-\frac{\D u(x)}{v}, \frac{1}{v}\right), \qquad v = \sqrt{1+|\D u|^2}, \ x \in \Omega, \ z \in \R.
\]

The Gauss map takes values in the upper hemisphere

\[
\Zset = \{X=(x_1,x_2,x_3)\in \R^3 \mid |X|=1, \ x_3>0\},
\]

and is
\begin{equation}
G(x,x_3)=\unitv(x,x_3).
\tag{16.86}
\end{equation}
In the normalized coordinates of Section 16.5,

\[
\psi(\zeta) = -\left(1 + |\D \uvec(\zeta)|^2\right)^{-1/2}\D \uvec(\zeta), \qquad \zeta \in \Uset,
\]

so (16.74)--(16.75) give

\[
J(G)(Y)=D_{11}\uvec(0)D_{22}\uvec(0)-(D_{12}\uvec(0))^2
\]



and

\[
|\partial G|^{2}(Y)=(D_{11}\tilde{u}(0))^{2}+2(D_{12}\tilde{u}(0))^{2}+(D_{22}\tilde{u}(0))^{2}.
\]

Equivalently,

\begin{equation}
\begin{aligned}
J(G)(Y)&=\kappa_{1}\kappa_{2},\\
|\partial G|^{2}(Y)&=\kappa_{1}^{2}+\kappa_{2}^{2}.
\end{aligned}
\tag{16.87}
\end{equation}

where $\kappa_1,\kappa_2$ are the principal curvatures.  Thus $G$ is
$(K,K')$-quasiconformal exactly when

\begin{equation}
\kappa_{1}^{2}+\kappa_{2}^{2}\leq 2K\kappa_{1}\kappa_{2}+K'.
\tag{16.88}
\end{equation}

Squaring (16.63') gives

\[
\frac{\alpha_{1}}{\alpha_{2}}\kappa_{1}^{2}+\frac{\alpha_{2}}{\alpha_{1}}\kappa_{2}^{2}=-2\kappa_{1}\kappa_{2}+\frac{\beta^{2}}{\alpha_{1}\alpha_{2}}.
\]

Hence (16.64')--(16.65') make the Gauss map of a solution of
(16.63)--(16.65) quasiconformal with

\begin{equation}
K=-\gamma,\qquad K'=\gamma\mu^{2}.
\tag{16.89}
\end{equation}

For Theorem 16.15, take

\[
\omega(X)=\left(-\frac{x_{2}}{1+x_{3}},\frac{x_{1}}{1+x_{3}},0\right).
\]

so $\Lambda_0=4$.  Lemma 16.13 and (16.87) give

\[
\int_{\mathfrak{S}_{R/2}(Y)}(\kappa_{1}^{2}+\kappa_{2}^{2})\,dA\leq C,
\]





where $C = C(K, K'R^2, \Lambda_1)$. Thus, since

\[
\kappa_1^2+\kappa_2^2 \geq \frac12(\kappa_1+\kappa_2)^2 = 2H^2,
\]

so one may take $\Lambda_2=C/2$.  The next lemma controls $\Lambda_1$.

\begin{lemma}[Area growth for quasiconformal Gauss maps]\label{gilbarg-trudinger.lem.16.16.area-growth}\dcref{ch16:16.16}\group{chapter-16}
Suppose $G$ is $(K,K')$-quasiconformal. Then

\begin{equation}
\tag{16.90}
|\Xset_{\rho/2}(Z)| \leq C\rho^2
\end{equation}

for each $Z\in\Sset$ and $\rho>0$ with $\Xset_\rho(Z)\subset\Xset_R(Y)$, where $C=C(K,K'R^2)$.
\end{lemma}

\begin{proof}
\sketch Testing the curvature identity (16.44) with a cutoff supported in $B_\rho$ and using (16.88) gives (16.91). Cauchy--Schwarz bounds the mean-curvature term by $C\rho$, and the graph area estimate (16.53) yields (16.90).


\begin{equation}
\tag{16.91}
\int_{\Omega} (\kappa_1^2 + \kappa_2^2)\eta^2\,dx \leq C
\end{equation}

The cutoff bounds and (16.53) give (16.90).
\end{proof}

Lemma 16.16 and Theorem 16.15 imply

\begin{equation}
\tag{16.92}
\sup_{x \in \mathfrak{S}_{\rho}^*(Y)} |\nu(X) - \nu(Y)| \leq C(\rho/R)^{\alpha}, \qquad \rho \in (0, R),
\end{equation}

where $C,\alpha>0$ depend only on $K,K'R^2$; for large $\rho$, enlarge
$C$ and use $|\nu|=1$.  Choose an isometry $P$ such that

\[
\vtilde(0) = P\nu(y, y_3) = (0, 0, 1)
\]

and let

\[
\Stilde = \{(\zeta, \zeta_3) \in \R^3 \mid (\zeta, \zeta_3) = P(x - y, x_3 - y_3),\ (x, x_3) \in \mathfrak{S}_{\theta R}^*(Y)\},
\]

For small $\theta$, the transformed component is a graph with

\begin{equation}
\tag{16.93}
\Stilde = \operatorname{graph}\tilde u = \{(\zeta, \zeta_3) \in \R^3 \mid \zeta \in \mathcal U,\, \zeta_3 = \tilde u(\zeta)\}.
\end{equation}

Furthermore, writing

\[
\vtilde(\zeta) = (1 + |D\tilde u(\zeta)|^2)^{-1/2}(-D\tilde u(\zeta), 1), \qquad \zeta \in \mathcal U,
\]






(16.92) gives

\[
|\vtilde(\zeta)-\vtilde(0)| \le C\theta^\alpha, \qquad \zeta\in\U,
\]

where $C,\alpha$ are as in (16.92). Consequently

\[
(1+|Du(\zeta)|^2)^{-1}|Du(\zeta)|^2+\bigl[(1+|Du(\zeta)|^2)^{-1/2}-1\bigr]^2\le (C\theta^\alpha)^2,\qquad \zeta\in\U.
\]

and therefore
\begin{equation}
|D\widetilde u(\zeta)|
\le[1-(C\theta^\alpha)^2]^{-1/2}C\theta^\alpha<\frac12,
\tag{16.94}
\end{equation}
provided
\begin{equation}
C\theta^\alpha<\frac13.
\tag{16.95}
\end{equation}
This slope bound continues the graphical representation and ensures
\begin{equation}
B_{\theta R/2}(0)\subset\mathcal U.
\tag{16.96}
\end{equation}

\begin{lemma}[Local connectivity of quasiconformal graphs]\label{gilbarg-trudinger.lem.16.17.local-connectivity}\dcref{ch16:16.17}\group{chapter-16}
Suppose $G$ is $(K,K')$-quasiconformal. There is $\theta\in(0,1)$, depending only on $K,K'R^2$, such that $\mathscr G_\rho(Y)$ is connected for every $\rho<\theta R$.
\end{lemma}

\begin{proof}
\sketch Choose $\rho=\theta R/2$ and let $\mathcal G_\beta$ be the components of
$\mathscr G_{\rho/2}(Y)$ meeting $B_{\beta\rho}$. The graphical representation
(16.93)--(16.96) gives
\begin{equation}\tag{16.97}
\dist(\pi_1,\pi_2)\le C_1(\beta+\theta^\alpha)\rho.
\end{equation}
If two adjacent components existed, the intervening region $\mathcal V$ would satisfy
\[
|\mathcal V|\le C_2(\beta+\theta^\alpha)\rho^3,
\]
\begin{equation}\tag{16.98}
A(\mathcal V\cap\partial\overline B_{\rho/2})
\le C_3(\beta+\theta^\alpha)\rho^2.
\end{equation}
The divergence theorem, (16.91), and Cauchy--Schwarz then imply
\begin{equation}\tag{16.99}
A(\mathcal G_1)+A(\mathcal G_2)
\le C_5\sqrt{\beta+\theta^\alpha}\,\rho^2.
\end{equation}
On the other hand, every component has
\begin{equation}\tag{16.100}
A(\mathcal G)\ge C_6\rho^2.
\end{equation}
For $\beta,\theta$ small enough this is a contradiction, so only the component
containing $Y$ remains and $\mathscr G_{\beta\rho}(Y)$ is connected.
\end{proof}

\begin{theorem}[H\"older estimate for the Gauss map]\label{gilbarg-trudinger.thm.16.18.gauss-map-holder}\dcref{ch16:16.18}\group{chapter-16}\uses{gilbarg-trudinger.lem.16.17.local-connectivity}\uses{gilbarg-trudinger.thm.16.15.quasiconformal-holder}
Suppose $G$ is $(K,K')$-quasiconformal. Then

\begin{equation}
\tag{16.101}
\sup_{X\in\mathscr{G}_{\rho}(Y)}|\nu(X)-\nu(Y)|
\le C(\rho/R)^\alpha,\qquad \rho\in(0,R),
\end{equation}

where $C,\alpha>0$ depend only on $K,K'R^2$.
\end{theorem}
\begin{proof}
\sketch Theorem 16.15 gives (16.92) on the component containing $Y$. Lemma 16.17 identifies that component with the full intersection for small radii; for larger radii the bound follows from $|\nu|=1$ after enlarging the constant.
\end{proof}

\begin{remark}[Consequences of the Gauss-map estimate]\label{gilbarg-trudinger.remark.16.101-consequences}\dcref{ch16:16.102}\group{chapter-16}
The estimate (16.101) implies

\begin{equation}
\tag{16.102}
|\nu(X)-\nu(\bar X)|\le 2^\alpha C(|X-\bar X|/R)^\alpha
\end{equation}

for all $X,\bar X\in\mathscr G_{R/4}(Y)$.  If $K'=0$ and
$\Omega=\mathbb R^2$, letting $R\to\infty$ shows that $\nu$ is constant.
\end{remark}

\begin{corollary}[Bernstein consequence for quasiconformal Gauss maps]\label{gilbarg-trudinger.cor.16.19.quasiconformal-gauss-map-bernstein}\dcref{ch16:16.19}\group{chapter-16}\uses{gilbarg-trudinger.thm.16.18.gauss-map-holder}
If $G$ is $(K,0)$-quasiconformal and $\Omega=\mathbb R^2$, then $u$ is affine.
\end{corollary}
\begin{proof}
\sketch Let $R\to\infty$ in (16.101). The Gauss map is constant, so $Du$ is constant.
\end{proof}

\section*{16.7. Applications to Equations of Mean Curvature Type}

Let $\mathscr G$ be the graph of a solution of (16.63)--(16.65). Its Gauss map is
$(K,K')$-quasiconformal with the constants in (16.89), so Corollary 16.19 yields
Theorem 16.12. For the curvature estimate below, extend $[a^{ij}]$ to a
$3\times3$ matrix by

\[
a^{3i}(x,z,p)=a^{i3}(x,z,p)=\sum_{j=1}^{2} p_j a^{ij}(x,z,p),\ i=1,2,
\qquad (x,z,p)\in\Omega\times\mathbb{R}\times\mathbb{R}^2,
\]

\[
a^{33}(x,z,p)=\sum_{i,j=1}^{2} p_i p_j a^{ij}(x,z,p).
\]

For $X=(x,z)$ and $v=(1+|p|^2)^{-1/2}(-p,1)$, set

\[
a_*^{ij}(X,v)=a^{ij}(x,z,p), \qquad i,j=1,2,3,
\]

\[
b_*(X,v)=(1+|p|^2)^{-1/2}b(x,z,p)
\]

on $(\Omega\times\mathbb R)\times\{\xi\in\mathbb R^3:|\xi|=1,\ \xi_3>0\}$.
Assume the H\"older conditions

\begin{equation}
\begin{aligned}
\left|a_*^{ij}(X,v)-a_*^{ij}(\bar X,\bar v)\right| &\le \mu_1\{(|X-\bar X|/R)^\beta+|v-\bar v|^\beta\},\\
\left|b_*(X,v)-b_*(\bar X,\bar v)\right| &\le \mu_2\{(|X-\bar X|/R)^\beta+|v-\bar v|^\beta\},
\end{aligned}
\tag{16.103}
\end{equation}

for all admissible $X,\bar X,v,\bar v$, where $\mu_1,\mu_2$ are constants and
$\beta\in(0,1)$.

\begin{theorem}[Interior curvature estimate]\label{gilbarg-trudinger.thm.16.20.interior-curvature-estimate}\dcref{ch16:16.20}\group{chapter-16}\uses{gilbarg-trudinger.thm.16.18.gauss-map-holder}\uses{gilbarg-trudinger.lem.16.17.local-connectivity}
Assume (16.63)--(16.65) and (16.103). If $\kappa_1,\kappa_2$ are the principal curvatures of $\mathcal S$ at $Y$, then

\begin{equation}
\kappa_1^2+\kappa_2^2\le C/R^2,
\tag{16.104}
\end{equation}

where $C=C(\gamma,\mu R,\mu_1,\mu_2R,\beta)$.
\end{theorem}

\begin{proof}
\sketch By Lemma 16.17, choose $\theta$ so that $\mathcal S_{\theta R}(Y)$ is
connected and has the representation (16.93)--(16.95). Applying the construction
of Section 16.4 to $\mathcal S$ and its transform gives

\begin{equation}
\tilde a^{ij}(\zeta)D_{ij}\tilde u+\tilde b(\zeta)=0
\tag{16.105}
\end{equation}



on $\mathcal U$, where

\[
|\xi|^2-\frac{(D\tilde u\cdot \xi)^2}{1+|D\tilde u|^2}
\leq \tilde a^{ij}\xi_i\xi_j
\leq \gamma\left[|\xi|^2-\frac{(D\tilde u\cdot \xi)^2}{1+|D\tilde u|^2}\right]
\]

for all $\xi\in \mathbb R^2$, and

\[
|\tilde b|\leq \mu\sqrt{1+|D\tilde u|^2}.
\]

Consequently, on $B=B_{R/2}(0)$, (16.94) gives

\begin{equation}\tag{16.106}
|\xi|^2/2\leq \tilde a^{ij}(\zeta)\xi_i\xi_j\leq \gamma|\xi|^2,
\qquad
|\tilde b(\zeta)|\leq 2\mu.
\end{equation}

Equations (16.102)--(16.103) and the coordinate change give

\begin{equation}\tag{16.107}
\left|\tilde a^{ij}(\zeta)-\tilde a^{ij}(\bar\zeta)\right|
\leq C(|\zeta-\bar\zeta|/R)^{\alpha\beta},
\qquad \zeta,\bar\zeta\in B,\ i,j=1,2
\end{equation}
\[
\left|\tilde b(\zeta)-\tilde b(\bar\zeta)\right|
\leq C(|\zeta-\bar\zeta|/R)^{\alpha\beta},
\qquad \zeta,\bar\zeta\in B.
\]

Since $\tilde u(0)=0$, (16.94) also gives

\begin{equation}\tag{16.108}
\sup_B |\tilde u|\leq R.
\end{equation}

Theorem 6.2 applied with (16.106)--(16.108) yields

\[
|\tilde u|_{2,\alpha\beta;B}^*\leq C\{R+\mu R^2+R^{2+\alpha\beta}\mu_2R^{-\alpha\beta}\}
\]
\[
\leq CR,
\]

where $C=C(\gamma,\mu R,\mu_1,\mu_2R,\beta)$. Hence

\[
\sum_{i,j=1}^2 |D_{ij}\tilde u(0)|\leq C/R,
\]

and (16.87) gives (16.104).
\end{proof}

The homogeneous equation requires only measurable coefficients.

\begin{theorem}[Interior gradient estimate in two variables]\label{gilbarg-trudinger.thm.16.21.two-dimensional-gradient-estimate}\dcref{ch16:16.21}\group{chapter-16}\uses{gilbarg-trudinger.lem.16.17.local-connectivity}
Assume (16.63), (16.64), $b\equiv0$, and measurability of $a^{ij}(x,u,Du)$ on $\Omega$.
With $m_R=\sup_{B_R(y)}(u-u(y))$,

\begin{equation}
|Du(y)| \le C_1 \exp(C_2 m_R / R).
\tag{16.109}
\end{equation}

where $C_1=C_1(\gamma)$ and $C_2=C_2(\gamma)$.
\end{theorem}

\begin{proof}
\sketch Choose $\theta=\theta(\gamma)$ as in Theorem 16.20 so that
$\Xi_{\theta R}(Y)$ is connected and (16.93)--(16.94) hold. Choose the matrix
$[p_{ij}]$ of $P$ with $p_{32}=0$. Then

\begin{equation}
(1+|Du(x)|^2)^{-1/2} = (1+|D\tilde u(\zeta)|^2)^{-1/2}(p_{31}D_1\tilde u(\zeta)-p_{33}).
\tag{16.110}
\end{equation}

Since $b\equiv0$, (16.105) becomes

\[
\alpha D_1 \tilde u + 2\beta D_{12}\tilde u + D_{22}\tilde u = 0
\]

where $\alpha=\tilde a_{11}/\tilde a_{22}$ and
$\beta=\tilde a_{12}/\tilde a_{22}$. Multiply by $p_{31}D_1\varphi$ for
$\varphi\in C_0^2(\mathcal U)$ and integrate by parts, using

\[
\int_{\mathcal U} D_{22}\tilde u\,D_1\varphi\,d\zeta
= -\int_{\mathcal U} D_2\tilde u\,D_{12}\varphi\,d\zeta
= \int_{\mathcal U} D_{21}\tilde u\,D_2\varphi\,d\zeta
\]

Set

\[
\psi = p_{31}D_1\tilde u - p_{33},
\]

Then

\[
\int_{\mathcal U} \left(\alpha D_1\psi\,D_1\varphi + 2\beta D_2\psi\,D_1\varphi + D_2\psi\,D_2\varphi\right)\,d\zeta = 0.
\]

so $\psi$ is a weak solution of the uniformly elliptic equation

\[
D_1(\alpha D_1\psi + 2\beta D_2\psi) + D_2(D_2\psi)=0
\qquad \text{on } \mathcal U.
\]

By (16.110), $\psi>0$ on $\mathcal U$. The Harnack inequality gives

\begin{equation}
\sup_{B_{R/2}(0)} \psi \le C \inf_{B_{R/2}(0)} \psi.
\tag{16.111}
\end{equation}



where $C=C(\gamma)$. Equations (16.94) and (16.110) imply

\[
(1+|Du(x)|^2)^{-1/2}\leq \psi(\zeta)\leq 2(1+|Du(x)|^2)^{-1/2}
\]

on $\mathcal U$. Define $v$ on $\Omega\times\mathbb R$ by

\[
v(x,z)=(1+|Du(x)|^2)^{1/2}, \qquad x\in \Omega, z\in \mathbb{R},
\]

Then (16.111) yields

\[
\sup_{X\in G_{R/2}} v(X)\leq C \inf_{X\in G_{R/2}} v(X),
\]

where $C=C(\gamma)$. Varying $Y$ gives $\chi=\chi(\gamma)\in(0,1/2)$ such that

\begin{equation}
\tag{16.112}
v(X_1)\leq Cv(X_2)
\end{equation}

whenever $X_i=(x^{(i)},u(x^{(i)}))$, $|X_1-X_2|\leq\chi R$, and
$x^{(i)}\in B_{R/2}(y)$. Let

\[
B_{\chi R}^{+}(y)=\{x\in B_{\chi R}\mid u(x)>u(y)\}
\]

and let $Y_1=(y^{(1)},u(y^{(1)}))$ and $Y_2=(y^{(2)},u(y^{(2)}))$ be such that $y^{(1)}, y^{(2)}\in \overline{B_{\chi R}^{+}(y)}$
and

\[
|Du(y^{(1)})|=\sup_{\overline{B_{\chi R}^{+}(y)}} |Du|,
\]

\[
|Du(y^{(2)})|=\inf_{\overline{B_{\chi R}^{+}(y)}} |Du|.
\]

Join $Y_1$ to $Y_2$ by points $X_1,\ldots,X_N$ in
$\mathcal G\cap(\overline{B_{\chi R}^{+}(y)}\times\mathbb R)$ with
$|X_{i+1}-X_i|<\chi R$. Repeated application of (16.112) gives

\[
v(Y_1)\leq C^N v(Y_2);
\]

that is

\[
\sup_{\overline{B_{\chi R}^{+}(y)}} \sqrt{1+|Du|^2}\leq C^N \inf_{\overline{B_{\chi R}^{+}(y)}} \sqrt{1+|Du|^2}.
\]

The chain can be chosen with

\[
N\leq C(m_R/R+1),
\]


where $C=C(\gamma)$. Hence
\[
\sup_{B^+_{\chi R}(y)} \sqrt{1+|Du|^2}\leq \{C_1 \exp(C_2 m_R/R)\} \inf_{B^+_{\chi R}(y)} \sqrt{1+|Du|^2},
\]
where $C_1=C_1(\gamma)$ and $C_2=C_2(\gamma)$. Finally,
\[
\inf_{B^+_{\chi R}(y)} |Du|\leq \chi^{-1} m_R/R,
\]
which gives (16.109).
\end{proof}

\section*{16.8. Appendix: Elliptic Parametric Functionals}

\begin{definition}[Parametric functional representation]\label{gilbarg-trudinger.def.16.113.parametric-functional}\dcref{ch16:16.113--16.115}\group{chapter-16}
For a bounded domain $\Omega\subset\mathbb R^2$ and a $C^1$ map
$\mathbf Y=(Y_1,Y_2,Y_3):\overline\Omega\to\mathbb R^3$, let
\begin{equation}
\tag{16.113}
I(\mathbf Y)=\int_\Omega G(x,\mathbf Y,D_1\mathbf Y,D_2\mathbf Y)\,dx,
\end{equation}
where $G$ is continuous on $\mathbb R^2\times\mathbb R^3\times\mathbb R^6$ and
$D_i\mathbf Y=(D_iY_1,D_iY_2,D_iY_3)$. Invariance under every
orientation-preserving diffeomorphism $\psi:\mathbb R^2\to\mathbb R^2$ means
\[
\int_{\Omega'} G(\zeta,\widetilde{\mathbf{Y}}(\zeta),D_1\widetilde{\mathbf{Y}}(\zeta),D_2\widetilde{\mathbf{Y}}(\zeta))\,d\zeta
=\int_\Omega G(x,\mathbf{Y}(x),D_1\mathbf{Y}(x),D_2\mathbf{Y}(x))\,dx,
\]
where $\Omega'=\psi(\Omega)$ and $\widetilde{\mathbf Y}=\mathbf Y\circ\psi^{-1}$.
This holds for all such $\psi$ and $\Omega$ if and only if there is a real-valued
$F$ on $\mathbb R^3\times\mathbb R^3$ such that
\begin{equation}
\tag{16.114}
G(x,X,p)=F(X,P), \qquad (x,X,p)\in \mathbb{R}^2\times \mathbb{R}^3\times \mathbb{R}^6,
\end{equation}
where
\[
P=(p_3p_5-p_2p_6,\; p_1p_6-p_3p_4,\; p_2p_4-p_1p_5),
\]






and
\begin{equation}
\tag{16.115}
F(X,\lambda q)=\lambda F(X,q), \qquad (X,q)\in\R^3\times\R^3,\quad \lambda>0.
\end{equation}
\end{definition}

Equation (16.114) makes $G$ independent of $x$. For
$p=(D_1\Y,D_2\Y)$,

\[
P=(D_1Y_3\cdot D_2Y_2-D_1Y_2\cdot D_2Y_3,\; D_1Y_1\cdot D_2Y_3-D_1Y_3\cdot D_2Y_1,
\]
\[
\hphantom{P=(}\; D_1Y_2\cdot D_2Y_1-D_1Y_1\cdot D_2Y_2).
\]

If $\Y$ is injective and $[D_jY_i]$ has rank two, then

\[
P=\chi\nu,
\]

where $\nu$ is the oriented unit normal of
$\G=\{Y(x):x\in\Omega\}$ and $\chi>0$ is the area magnification factor. Thus

\[
I(\Y)=\int_{\G} F(X,\nu(X))\,dA(X);
\]

so $I$ depends only on the oriented surface.

\begin{definition}[Elliptic parametric functional]\label{gilbarg-trudinger.def.16.116.elliptic-parametric-functional}\dcref{ch16:16.116--16.117}\group{chapter-16}
For a smooth oriented surface $\G\subset\R^3$ of finite area, define
\begin{equation}
\tag{16.116}
J(\G)=\int_{\G}F(X,\nu(X))\,dA(X).
\end{equation}
If $F$ satisfies (16.115), this is a parametric functional and
$J(\G)=I(\Y)$ for every injective rank-two parametrization $\Y$ of $\G$.
It is elliptic if $F\in C^2(\R^3\times(\R^3-\{0\}))$ and
\begin{equation}
\tag{16.117}
D_{q_i q_j}F(X,q)\,\xi_i\xi_j\ge |q|^{-1}|\xi'|^2,\qquad
\xi'=\xi-\left(\xi\cdot\frac{q}{|q|}\right)\frac{q}{|q|}
\end{equation}
for all $X\in\R^3$, $q\ne0$, and $\xi\in\R^3$.
\end{definition}



For the graph

\[
\mathfrak{S}=\{(x,u(x))\in \mathbb{R}^{3}\mid x\in \Omega\},
\]

with $u\in C^2(\bar\Omega)$ and downward unit normal
\[
(Du,-1)/\sqrt{1+|Du|^{2}},
\]
we have

\[
J(\mathfrak{S})=\int_{\Omega}F(x,u(x),Du(x),-1)\,dx.
\]

using $dA=\sqrt{1+|Du|^2}\,dx$. Its Euler--Lagrange equation is

\[
\sum_{i=1}^{2}D_{i}[D_{q_{i}}F(x,u,Du,-1)]-D_{x_{3}}F(x,u,Du,-1)=0.
\]

The chain rule and (16.115) put this equation in the form

\[
Qu=a^{ij}(x,u,Du)D_{i,j}u+b(x,u,Du)=0,
\]

where

\begin{equation}
\tag{16.118}
a^{ij}(x,u,Du)=\nu D_{q_{i}q_{j}}F(x,u,Du,-1),\qquad i,j=1,2,
\end{equation}

\begin{equation}
\tag{16.119}
b(x,u,Du)=\nu\sum_{i=1}^{3}D_{q_{i}x_{i}}F(x,u,Du,-1),
\end{equation}

where $\nu=\sqrt{1+|Du|^2}$. Conditions (16.115) and (16.117) imply
(16.64)--(16.65), with $\gamma$ and $\mu$ determined by $F$; hence this is an
equation of mean-curvature type.

For (16.118)--(16.119), the coefficients introduced in Section 16.7 are

\[
a^{ij}_{*}(X,\nu)=D_{q_{i}q_{j}}F(X,\nu),\qquad i,j=1,2,3.
\]

\[
b_{*}(X,\nu)=\sum_{i=1}^{3}D_{q_{i}x_{i}}F(X,\nu);
\]

If $F\in C^3(\mathbb R^3\times(\mathbb R^3-\{0\}))$, these coefficients satisfy
(16.103), with $\mu_1,\mu_2$ determined by $F$.
